% Figure from MATH 493 notes, section 6. Compile with the notes preamble.
\begin{tikzpicture}[
    num/.style={circle, fill=white, inner sep=0.6pt, font=\scriptsize}]
  % the integers 0..35 wound onto a spiral, 12 per turn, clockwise from the top
  \draw[black!30, thin, domain=0:35, samples=180, smooth, variable=\t]
    plot ({(1.25+0.62*\t/12)*cos(90-30*\t)}, {(1.25+0.62*\t/12)*sin(90-30*\t)});
  % the residue class [2] = {2, 14, 26, ...}: one ray
  \draw[figR, thick] (30:0.45) -- (30:3.1);
  \node[font=\small, figR, anchor=south west, inner sep=1pt] at (30:3.1) {$[2]$};
  % the walk 9 -> 14 (add 5)
  \draw[figB, very thick, -{Stealth[length=5pt]}, domain=9:13.72, samples=60, smooth, variable=\t]
    plot ({(1.25+0.62*\t/12)*cos(90-30*\t)}, {(1.25+0.62*\t/12)*sin(90-30*\t)});
  \foreach \k in {0,...,35} {
    \pgfmathsetmacro{\r}{1.25+0.62*\k/12}
    \pgfmathtruncatemacro{\c}{mod(\k,12)==2 ? 1 : (\k==9 ? 2 : 0)}
    \ifcase\c
      \node[num] at ({90-30*\k}:\r) {$\k$};
    \or
      \node[num, text=figR, draw=figR, thin] at ({90-30*\k}:\r) {$\k$};
    \or
      \node[num, text=figB, draw=figB, thin] at ({90-30*\k}:\r) {$\k$};
    \fi
  }
  \node[font=\scriptsize, figB] at (104:1.55) {$+5$};
\end{tikzpicture}
